% =====================================================================
%  Annexes
% =====================================================================

\chapter{Formulaire récapitulatif}
\label{ann:formulaire}

Cette annexe rassemble les fiches \og tableau blanc\fg{} de tous les chapitres. Elle est générée automatiquement par \texttt{docs/cours/make\_formulaire.py}.

\section*{Chapitre 1 — Arbitrage, stationnarité et retour à la moyenne}
\addcontentsline{toc}{section}{Chapitre 1 — Arbitrage, stationnarité et retour à la moyenne}
\begin{itemize}[leftmargin=1.2em,itemsep=0.05em]
\item \textbf{Arbitrage pur :} coût $\le 0$, perte impossible, gain possible. \textbf{Statistique :} $v(0)=0$, $\E[v(t)]\to>0$, $\Prob(v(t)<0)\to0$, $\Var(v(t))/t\to0$.
\item \textbf{Variance d'un écart :} $\Var(Y-\beta X) = \sigma_Y^2 + \beta^2\sigma_X^2 - 2\beta\Cov(X,Y)$.
\item \textbf{Couverture optimale :} $\beta^\star = \Cov(X,Y)/\sigma_X^2 = \rho\,\sigma_Y/\sigma_X$ ; variance restante $\sigma_Y^2(1-\rho^2)$.
\item \textbf{Stationnarité faible :} $\E X_t = \mu$, $\Var X_t = \gamma(0)$, $\Cov(X_t,X_{t-h}) = \gamma(h)$, indépendants de $t$.
\item \textbf{Marche aléatoire :} $P_t = P_0 + \sum\varepsilon_i$ ; $\E P_t = P_0$ ; $\Var P_t = t\sigma^2$ ; non stationnaire ; $I(1)$.
\item \textbf{AR(1) :} $X_t - \mu = \varphi(X_{t-1}-\mu) + \varepsilon_t$ ; stationnaire si $|\varphi|<1$ ; $\mu = c/(1-\varphi)$ ; $\gamma(0) = \sigma^2/(1-\varphi^2)$ ; $\rho(h) = \varphi^h$.
\item \textbf{Retour à la moyenne :} $\Delta X_t = -(1-\varphi)(X_{t-1}-\mu) + \varepsilon_t$.
\item \textbf{Demi-vie :} $h_{1/2} = -\ln2/\ln\varphi \approx \ln 2/(1-\varphi)$.
\item \textbf{Cointégration (intuition) :} $X, Y \sim I(1)$ et $Y - \beta X \sim I(0)$.
\item \textbf{Z-score :} $z_t = (S_t - m)/d$, avec $m$, $d$ estimés \emph{sur la formation seulement}.
\item \textbf{Gain sans anticipation :} $G_{t+1} = q_t(S_{t+1}-S_t)$.
\end{itemize}

\section*{Chapitre 2 — Régression et estimation}
\addcontentsline{toc}{section}{Chapitre 2 — Régression et estimation}
\begin{itemize}[leftmargin=1.2em,itemsep=0.05em]
\item \textbf{MCO simples :} $\hat b = \sum(x-\bar x)(y-\bar y)/\sum(x-\bar x)^2$, $\hat a = \bar y - \hat b\bar x$.
\item \textbf{MCO matriciels :} $\hat\beta = (X^\top X)^{-1}X^\top y$ ; $X^\top e = 0$ ; $\hat\sigma^2 = e^\top e/(n-k)$.
\item \textbf{Sans biais :} $\hat\beta = \beta + (X^\top X)^{-1}X^\top\varepsilon$ ; $\Var(\hat\beta) = \sigma^2(X^\top X)^{-1}$ sous (H4).
\item \textbf{Sandwich :} $(X^\top X)^{-1}\,\hat\Omega\,(X^\top X)^{-1}$ ; Newey-West : poids $1-j/(L+1)$.
\item \textbf{Asymétrie :} $\hat\beta_{Y|X}\hat\beta_{X|Y} = \hat\rho^2$.
\item \textbf{Orthogonale :} $s_{XY}b^2 + (s_{XX}-s_{YY})b - s_{XY} = 0$, racine de même signe que $s_{XY}$.
\item \textbf{Log-vraisemblance AR(1) :} $\ell = -\frac n2\ln(2\pi\sigma^2) - \frac1{2\sigma^2}\sum(x_t - c - \varphi x_{t-1})^2$ ; MV conditionnel $=$ MCO.
\item \textbf{Asymptotique :} $\Var(\hat\varphi)\approx(1-\varphi^2)/n$.
\item \textbf{Correction de Kendall :} $\hat\varphi_c = \hat\varphi + (1+3\hat\varphi)/n$.
\end{itemize}

\section*{Chapitre 3 — Tester la stationnarité et le retour à la moyenne}
\addcontentsline{toc}{section}{Chapitre 3 — Tester la stationnarité et le retour à la moyenne}
\begin{itemize}[leftmargin=1.2em,itemsep=0.05em]
\item \textbf{DF :} $\Delta x_t = c + \gamma x_{t-1} + \varepsilon_t$, $H_0 : \gamma = 0$ (racine unitaire), $\tau = \hat\gamma/\mathrm{se}$.
\item \textbf{Pourquoi non normale :} $\sum x_{t-1}\varepsilon_t = \frac12(x_n^2 - \sum\varepsilon_t^2)$ $\Rightarrow$ $\frac12(W^2-1)$, négatif 68\,\% du temps.
\item \textbf{Seuils 5\,\% :} $-1{,}95$ (rien), $-2{,}86$ (constante), $-3{,}41$ (tendance).
\item \textbf{ADF :} ajouter $\sum\delta_j\Delta x_{t-j}$, choisir $p$ par AIC.
\item \textbf{Puissance :} gouvernée par $n(1-\varphi)\approx 0{,}69\,n/h_{1/2}$ ; viser $\ge 30$ demi-vies.
\item \textbf{KPSS :} $H_0$ stationnarité ; $\eta = n^{-2}\sum S_t^2/\hat\sigma^2_{LR}$ ; seuil $0{,}463$.
\item \textbf{Ratio de variance :} $VR(q) = \Var(x_{t+q}-x_t)/(q\Var(\Delta x))$ ; AR(1) : $(1-\varphi^q)/(q(1-\varphi))$.
\item \textbf{Hurst :} $\mathrm{sd}(x_{t+\tau}-x_t)\propto\tau^H$ ; $H<\frac12$ retour, $=\frac12$ marche, $>\frac12$ tendance.
\end{itemize}

\section*{Chapitre 4 — Cointégration}
\addcontentsline{toc}{section}{Chapitre 4 — Cointégration}
\begin{itemize}[leftmargin=1.2em,itemsep=0.05em]
\item \textbf{$I(d)$ :} différencier $d$ fois pour être stationnaire. Prix $\sim I(1)$.
\item \textbf{Cointégration :} $x\sim I(1)$, $\beta^\top x\sim I(0)$ ; rang $r$ = nombre de relations indépendantes.
\item \textbf{Engle-Granger :} MCO en niveaux, puis ADF sur résidus ; seuils 5\,\% : $-3{,}34$ (2 var.), $-3{,}74$ (3 var.).
\item \textbf{Superconvergence :} $\hat b - b = O(1/n)$ si la paire est cointégrée.
\item \textbf{Correction d'erreur :} $\Delta x_{1t} = \alpha_1S_{t-1}+\dots$, $\Delta x_{2t} = \alpha_2S_{t-1}+\dots$ ; $\varphi_S = 1+\alpha_1-\beta\alpha_2$.
\item \textbf{VECM :} $\Delta x_t = \Pi x_{t-1} + \sum\Gamma_i\Delta x_{t-i} + \varepsilon_t$, $\Pi = \alpha\beta^\top$, rang $r$.
\item \textbf{Johansen :} trace $= -T\sum_{i>r}\ln(1-\hat\lambda_i)$ ; tester $r = 0, 1, \dots$
\item \textbf{Cointégration partielle :} $S = M + R$ ; $R^2_{MR} = 2\sigma_M^2/(2\sigma_M^2 + (1+\rho)\sigma_R^2)$.
\end{itemize}

\section*{Chapitre 5 — Le processus d'Ornstein-Uhlenbeck}
\addcontentsline{toc}{section}{Chapitre 5 — Le processus d'Ornstein-Uhlenbeck}
\begin{itemize}[leftmargin=1.2em,itemsep=0.05em]
\item \textbf{Brownien :} $W_t - W_s\sim\N(0,t-s)$, $(\dd W)^2 = \dd t$.
\item \textbf{Itô :} $\dd f = (\partial_tf + b\,\partial_xf + \frac12s^2\partial_{xx}f)\dd t + s\,\partial_xf\,\dd W$.
\item \textbf{OU :} $\dd X = \theta(\mu - X)\dd t + \sigma\dd W$ ; $X_t = \mu + (X_0-\mu)e^{-\theta t} + \sigma\int_0^te^{-\theta(t-u)}\dd W_u$.
\item \textbf{Moments :} moyenne $\mu + (X_0-\mu)e^{-\theta t}$ ; variance $\frac{\sigma^2}{2\theta}(1-e^{-2\theta t})$ ; $\sigma_{eq} = \sigma/\sqrt{2\theta}$ ; $h_{1/2} = \ln2/\theta$.
\item \textbf{OU $\leftrightarrow$ AR(1) :} $\varphi = e^{-\theta\Delta}$, $\sigma_\eta^2 = \frac{\sigma^2}{2\theta}(1-\varphi^2)$ ; $\hat\theta = -\ln\hat\varphi/\Delta$.
\item \textbf{Kolmogorov :} $\frac12s^2u'' + bu' = -1$.
\item \textbf{Retour :} $T_{a\to0} = \frac2{\sigma^2}\int_0^ae^{ky^2}\int_y^\infty e^{-kz^2}\dd z\,\dd y$, $k=\theta/\sigma^2$.
\item \textbf{Seuil optimal :} maximiser $(a-c)/(T_{0\to\pm a}+T_{a\to0})$ ; le seuil monte avec les coûts.
\end{itemize}

\section*{Chapitre 6 — La méthode des distances}
\addcontentsline{toc}{section}{Chapitre 6 — La méthode des distances}
\begin{itemize}[leftmargin=1.2em,itemsep=0.05em]
\item \textbf{Formation 12 mois :} $\tilde P_{i,t} = P_{i,t}/P_{i,0}$, $\mathrm{SSD}_{ij} = \sum_t(\tilde P_{i,t}-\tilde P_{j,t})^2$, top 20 paires.
\item \textbf{Calcul matriciel :} $\mathrm{SSD} = \mathrm{diag}(G)\mathbf 1^\top + \mathbf 1\mathrm{diag}(G)^\top - 2G$, $G = \tilde P^\top\tilde P$.
\item \textbf{Trading 6 mois :} ouvrir à $|D_t| > 2\hat\sigma$ (1\,\$ long / 1\,\$ court), fermer au croisement.
\item \textbf{Gain d'un trade :} $(P^L_{t_1}/P^L_{t_0}-1) - (P^C_{t_1}/P^C_{t_0}-1)$.
\item \textbf{Capital engagé :} somme des gains / 20 ; \textbf{employé :} / nombre de paires ouvertes.
\item \textbf{Biais de la SSD :} $\mathrm{SSD}/T = \bar D^2 + \widehat\Var(D)$ : favorise les écarts peu volatils.
\end{itemize}

\section*{Chapitre 7 — Cointégration appliquée : sélectionner les paires sans se mentir}
\addcontentsline{toc}{section}{Chapitre 7 — Cointégration appliquée : sélectionner les paires sans se mentir}
\begin{itemize}[leftmargin=1.2em,itemsep=0.05em]
\item \textbf{Faux positifs attendus :} $m\alpha$, avec $m = N(N-1)/2$.
\item \textbf{FWER} $=\Prob(V\ge1)$ ; \textbf{FDR} $=\E[V/\max(R,1)]$.
\item \textbf{Bonferroni :} rejeter si $p < \alpha/m$ ; preuve par l'inégalité de la réunion.
\item \textbf{BH :} plus grand $k$ tel que $p_{(k)}\le kq/m$ ; rejeter les $k$ plus petites.
\item \textbf{Pipeline :} a priori économique $\to$ test corrigé $\to$ demi-vie et amplitude $\to$ stabilité $\to$ contraintes pratiques.
\item \textbf{Erreur-type honnête :} regrouper par univers ou période, jamais par paire si les paires partagent des actifs.
\item \textbf{Taille :} notionnel $K = R/\sigma_S$.
\end{itemize}

\section*{Chapitre 8 — Modèles espace-état et filtre de Kalman}
\addcontentsline{toc}{section}{Chapitre 8 — Modèles espace-état et filtre de Kalman}
\begin{itemize}[leftmargin=1.2em,itemsep=0.05em]
\item \textbf{Modèle :} $\xi_t = F\xi_{t-1} + w_t$ ($Q$) ; $y_t = H_t\xi_t + v_t$ ($R$).
\item \textbf{Conditionnement :} $\E[X|y] = \mu_X + \Sigma_{XY}\Sigma_{YY}^{-1}(y-\mu_Y)$ ; $\Var = \Sigma_{XX} - \Sigma_{XY}\Sigma_{YY}^{-1}\Sigma_{YX}$.
\item \textbf{Prédiction :} $m^- = Fm$, $P^- = FPF^\top + Q$.
\item \textbf{Innovation :} $e = y - Hm^-$, $F_t = HP^-H^\top + R$.
\item \textbf{Mise à jour :} $K = P^-H^\top F_t^{-1}$, $m = m^- + Ke$, $P = P^- - KHP^-$.
\item \textbf{Scalaire :} $K = P/(P+R)$.
\item \textbf{Vraisemblance :} $\ell = -\frac12\sum(\ln 2\pi F_t + e_t^2/F_t)$.
\item \textbf{Ratio variable :} $\xi = (\alpha,\beta)$, $F = I$, $Q = \frac\delta{1-\delta}I$, $H_t = (1,x_t)$ ; signal $e_t/\sqrt{F_t}$.
\end{itemize}

\section*{Chapitre 9 — Copules}
\addcontentsline{toc}{section}{Chapitre 9 — Copules}
\begin{itemize}[leftmargin=1.2em,itemsep=0.05em]
\item \textbf{Transformation :} $U = F_X(X)\sim\mathcal U[0,1]$.
\item \textbf{Sklar :} $H(x,y) = C(F_X(x),F_Y(y))$.
\item \textbf{Kendall :} $\tau = \Prob(\text{concordance}) - \Prob(\text{discordance})$ ; gaussienne/Student : $\rho = \sin(\pi\tau/2)$ ; Clayton : $\theta = 2\tau/(1-\tau)$.
\item \textbf{Queue :} $\lambda_L = \lim_{q\to0}\Prob(V\le q\mid U\le q)$ ; Clayton $2^{-1/\theta}$ ; gaussienne 0.
\item \textbf{Clayton :} $C = (u^{-\theta}+v^{-\theta}-1)^{-1/\theta}$.
\item \textbf{Fonction $h$ :} $h(u|v) = \partial C/\partial v$ ; gaussienne $\Phi\big((\Phi^{-1}(u)-\rho\Phi^{-1}(v))/\sqrt{1-\rho^2}\big)$.
\item \textbf{Pseudo-observations :} rang$/(n+1)$ ; choix par AIC $= 2k-2\ell$.
\end{itemize}

\section*{Chapitre 10 — Arbitrage statistique multi-actifs}
\addcontentsline{toc}{section}{Chapitre 10 — Arbitrage statistique multi-actifs}
\begin{itemize}[leftmargin=1.2em,itemsep=0.05em]
\item \textbf{ACP :} $\max_w w^\top Cw$ s.c. $w^\top w=1$ $\Rightarrow$ $Cw = \lambda w$ ; part expliquée $\sum_{j\le k}\lambda_j/N$.
\item \textbf{Portefeuille propre :} $Q_{ij} = v_{j,i}/\sigma_i$, $F_{jt} = \sum_iQ_{ij}R_{it}$.
\item \textbf{Marchenko-Pastur :} bruit entre $(1\pm\sqrt{N/T})^2$.
\item \textbf{Modèle :} $R_i = \alpha_i + \sum_j\beta_{ij}F_j + \dd X_i$, $X_i$ OU.
\item \textbf{AR(1) sur $X$ :} $\kappa = -252\ln b$, $m = a/(1-b)$, $\sigma_{eq} = \sqrt{\Var(\zeta)/(1-b^2)}$.
\item \textbf{s-score :} $s = (X - m)/\sigma_{eq}$ ; ouvrir à $\pm1{,}25$, fermer court à $0{,}75$, long à $-0{,}5$ ; filtre $\kappa > 252/30$.
\item \textbf{Box-Tiao :} minimiser $w^\top A\Gamma A^\top w/w^\top\Gamma w$.
\end{itemize}

\section*{Chapitre 11 — Apprentissage automatique}
\addcontentsline{toc}{section}{Chapitre 11 — Apprentissage automatique}
\begin{itemize}[leftmargin=1.2em,itemsep=0.05em]
\item \textbf{Variable :} calculable en $t$ ; \textbf{étiquette :} regarde vers $t+H$.
\item \textbf{Chevauchement :} étiquettes à horizon $H$ partagent $H-1$ jours $\Rightarrow$ K-fold mélangé fuit.
\item \textbf{Remède :} walk-forward, purge des dates qui chevauchent le test, embargo.
\item \textbf{Test de base :} toujours comparer à une règle simple sur le même échantillon de test.
\item \textbf{Sélection :} ACP (expositions) $\to$ OPTICS (groupes) $\to$ cointégration intra-groupe $+$ BH.
\end{itemize}

\section*{Chapitre 12 — Microstructure et coûts de transaction}
\addcontentsline{toc}{section}{Chapitre 12 — Microstructure et coûts de transaction}
\begin{itemize}[leftmargin=1.2em,itemsep=0.05em]
\item \textbf{Fourchette :} ask $-$ bid ; mid $=$ (bid+ask)/2 ; taker paie la demi-fourchette, maker la reçoit mais subit la sélection adverse.
\item \textbf{Roll :} $\Cov(\Delta p_t,\Delta p_{t-1}) = -s^2/4$ $\Rightarrow$ $s = 2\sqrt{-\Cov}$ (données de transaction seulement).
\item \textbf{Aller-retour de paire :} $4c$ + emprunt.
\item \textbf{Règle :} $\sigma_{eq}\ge$ coût d'un aller-retour.
\item \textbf{Impact :} $Y\sigma_{\text{jour}}\sqrt{Q/V}$.
\item \textbf{Backtest :} exécuter au plus tôt en $t+1$ ; payer la demi-fourchette.
\end{itemize}

\section*{Chapitre 13 — Arbitrage spatial et triangulaire}
\addcontentsline{toc}{section}{Chapitre 13 — Arbitrage spatial et triangulaire}
\begin{itemize}[leftmargin=1.2em,itemsep=0.05em]
\item \textbf{Spatial :} opportunité si $b_B/a_A > (1+f_A)/(1-f_B)\approx1+f_A+f_B$.
\item \textbf{Inventaire :} rééquilibrage rentable si $Q > 10^4K/e$ ($K$ coût de transfert, $e$ avantage en pb).
\item \textbf{Écart persistant :} chercher la friction (contrôle des capitaux, transfert, contrepartie).
\item \textbf{Latence :} le premier prend tout ; durée de vie en millisecondes $\Rightarrow$ inaccessible au particulier.
\item \textbf{Triangulaire :} profitable si $z^{bid}(1-f)^3/(x^{ask}y^{ask}) > 1$ ; probabilité $2\Prob(Z > 3(h+f)/(\sigma_n\sqrt3))$.
\item \textbf{Ligne rouge :} arbitrage de latence contre un courtier ou une \emph{prop firm} : interdit.
\end{itemize}

\section*{Chapitre 14 — Arbitrage de base, \emph{cash-and-carry} et \emph{funding}}
\addcontentsline{toc}{section}{Chapitre 14 — Arbitrage de base, \emph{cash-and-carry} et \emph{funding}}
\begin{itemize}[leftmargin=1.2em,itemsep=0.05em]
\item \textbf{Portage :} $F = Se^{(r-q)T}$ ; base $F - S\to0$ à l'échéance.
\item \textbf{Bande :} $Se^{(r_l-q)T} - c' \le F \le Se^{(r_b-q)T} + c$.
\item \textbf{Cash-and-carry :} long comptant financé, court future ; gain final fixé, besoin de marge variable.
\item \textbf{Funding :} longs paient les courts si perpétuel $>$ indice ; formule propre à chaque plateforme.
\item \textbf{Liquidation :} hausse de $\approx 1/L - m$ sur la jambe courte.
\item \textbf{atjte :} $s_t = P^{\text{plat}} - kP^{\text{MT5}}$ ; maker sur la plateforme, taker sur MT5 ; centre théorique $=$ référence $\times$ portage $\times$ jours$/365$.
\end{itemize}

\section*{Chapitre 15 — Exécution, couverture et architecture logicielle}
\addcontentsline{toc}{section}{Chapitre 15 — Exécution, couverture et architecture logicielle}
\begin{itemize}[leftmargin=1.2em,itemsep=0.05em]
\item \textbf{Risque de jambe :} écart-type $\sigma\sqrt\Delta$ ; nul en espérance seulement si le délai est indépendant des prix.
\item \textbf{Sélection adverse :} un ordre maker est exécuté quand le marché lui donne tort ; coût croissant avec le retard de recotation.
\item \textbf{Où coter :} sur la plateforme suiveuse ; couvrir sur la meneuse.
\item \textbf{Architecture :} passerelles, contrôle du ratio, lot minimal, réconciliation sur positions réelles, sommeil si flux périmé, un robot par contrat, mode simulation, journaux.
\end{itemize}

\section*{Chapitre 16 — \emph{Backtest} rigoureux}
\addcontentsline{toc}{section}{Chapitre 16 — \emph{Backtest} rigoureux}
\begin{itemize}[leftmargin=1.2em,itemsep=0.05em]
\item \textbf{Sans anticipation :} $\text{gain}_{t+1} = q_t(S_{t+1}-S_t)$, $q_t$ fonction des seules données $\le t$ ; test automatique : modifier le futur ne change pas le passé.
\item \textbf{Trois biais :} anticipation, survie (garder les radiés), sélection (compter les essais).
\item \textbf{Erreur-type du Sharpe :} $\sqrt{(1+SR^2/2)/T}$ par période ; annualisé $\approx 1/\sqrt{\text{années}}$.
\item \textbf{Maximum de $N$ Sharpe nuls :} $\approx\sqrt V\big((1-\gamma)\Phi^{-1}(1-\tfrac1N) + \gamma\Phi^{-1}(1-\tfrac{1}{Ne})\big)$, croît comme $\sqrt{2\ln N/T}$.
\item \textbf{PSR :} $\Phi\big((\widehat{SR}-SR^\ast)\sqrt{T-1}/\sqrt{1-\gamma_3\widehat{SR}+\frac{\gamma_4-1}{4}\widehat{SR}^2}\big)$ ; \textbf{DSR :} PSR avec $SR^\ast = \E[\max]$.
\item \textbf{PBO :} proportion des découpages CSCV où la meilleure variante en sélection finit sous la médiane en test.
\item \textbf{Découpage :} formation, validation, test regardé une seule fois ; \emph{walk-forward} ; purge et embargo.
\end{itemize}

\section*{Chapitre 17 — Risque et limites de l'arbitrage}
\addcontentsline{toc}{section}{Chapitre 17 — Risque et limites de l'arbitrage}
\begin{itemize}[leftmargin=1.2em,itemsep=0.05em]
\item \textbf{Kelly :} $g(f) \approx f\mu - \tfrac12 f^2\sigma^2$ ; $f^\star = \mu/\sigma^2$ ; $g(f^\star) = SR^2/2$ ; $g(2f^\star) = 0$ ; en pratique $0{,}1$ à $0{,}25\,f^\star$.
\item \textbf{Ciblage de volatilité :} $w_t = \min(\sigma^\ast/\hat\sigma_t, L_{\max})$, $\hat\sigma_t$ estimée sur le passé.
\item \textbf{Protections :} stop $|z|>4$ (quasi gratuit) ; durée $> 3$ demi-vies ; CUSUM $S^\pm_t = \max(0, S^\pm_{t-1} \pm z_t - k)$, alarme si $> h$.
\item \textbf{Limites de l'arbitrage :} perte latente $\Rightarrow$ appel de marge $\Rightarrow$ liquidation au pire moment ; levier choisi sur le pire élargissement.
\item \textbf{Diversification :} $\sigma_P = \sigma\sqrt{1/N + \rho(1-1/N)}$ ; $N_{\text{eff}} = N/(1+(N-1)\rho)$ ; $\rho$ monte en crise.
\item \textbf{Hors statistique :} plateforme, contrepartie, courtier, emprunt de titres, opérations sur titres, réglementation, \emph{prop firms}, erreurs opérationnelles.
\end{itemize}

\section*{Chapitre 18 — Du papier au code avec Claude Code}
\addcontentsline{toc}{section}{Chapitre 18 — Du papier au code avec Claude Code}
\begin{itemize}[leftmargin=1.2em,itemsep=0.05em]
\item \textbf{Cycle :} spécification $\to$ tests $\to$ code $\to$ recherche $\to$ critères $\to$ \emph{paper} $\to$ petite taille $\to$ taille cible $\to$ revue.
\item \textbf{Parité :} version vectorisée = version événementielle, vérifiée par un test.
\item \textbf{Grille :} G0 hypothèse ; G1 $p_{EG}<0{,}05$ ; G2 demi-vie $\in[1,30]$ ; G3 DSR $>0{,}95$ ; G4 test net $>0$ ; G5 rentabilité nulle $> 2\times$ coûts.
\item \textbf{Retard :} un jour de retard coûte d'autant plus que la demi-vie est courte.
\item \textbf{Suivi :} $Z_T = (\widehat{SR}_{\text{réel}} - SR_{\text{bt}})\sqrt{T-1}$ ; contrôles datés à l'avance, seuil de Bonferroni.
\item \textbf{Journal :} coût d'implémentation = P\&L réel $-$ P\&L du \emph{backtest} rejoué.
\end{itemize}


% =====================================================================
\chapter{Notations}
\label{ann:notations}
% =====================================================================

\begin{center}\small
\begin{longtable}{>{\raggedright}p{3.2cm}>{\raggedright\arraybackslash}p{10.6cm}}
\toprule
Symbole & Signification \\
\midrule
\endhead
$t$, $T$, $n$ & date (en périodes), nombre total de périodes, taille d'échantillon \\
$X_t$, $Y_t$ ou $x_t$, $y_t$ & prix (ou log-prix) des deux jambes d'une paire \\
$S_t$, $s_t$ & écart (\emph{spread}) : $S_t = Y_t - \beta X_t$ (ou en log-prix) \\
$\beta$, $\hat\beta$ & ratio de couverture, et son estimation \\
$\varepsilon_t$, $\eta_t$ & bruit blanc (chocs non corrélés, moyenne nulle) \\
$\mu$, $\sigma$, $\sigma^2$ & moyenne, écart-type, variance \\
$\rho$, $\rho(h)$ & corrélation ; autocorrélation au retard $h$ \\
$\gamma(h)$ & autocovariance au retard $h$ \\
$\varphi$ & coefficient autorégressif de l'AR(1) \\
$h_{1/2}$ & demi-vie : $-\ln2/\ln\varphi$ (discret), $\ln2/\theta$ (continu) \\
$\theta$, $\sigma_{eq}$ & vitesse de retour et écart-type d'équilibre d'un processus d'Ornstein-Uhlenbeck \\
$W_t$ & mouvement brownien standard \\
$I(d)$ & intégré d'ordre $d$ (stationnaire après $d$ différenciations) \\
$\Delta x_t$ & différence première $x_t - x_{t-1}$ \\
$z_t$ & \emph{z-score} : $(S_t - m)/d$ avec $m$, $d$ estimés sur le passé \\
$q_t$ & position décidée à la clôture de $t$ (en unités d'écart), appliquée de $t$ à $t+1$ \\
$c$, pb & coût par jambe ; point de base $=10^{-4}$ \\
$\Phi$, $\Phi^{-1}$ & fonction de répartition et quantile de la loi normale centrée réduite \\
$SR$, $\widehat{SR}$ & ratio de Sharpe (par période ou annualisé, précisé à chaque usage) \\
PSR, DSR, PBO & Sharpe probabiliste, Sharpe déflaté, probabilité de surapprentissage \\
$f^\star$ & levier de Kelly $\mu/\sigma^2$ \\
$F$, $S$, $r$, $q$ & prix à terme, prix comptant, taux sans risque, rendement de détention (chapitre~14) \\
$\N(\mu,\sigma^2)$, $\iid$ & loi normale ; indépendantes et identiquement distribuées \\
\bottomrule
\end{longtable}
\end{center}

% =====================================================================
\chapter{Glossaire}
\label{ann:glossaire}
% =====================================================================

\begin{description}[leftmargin=0pt,itemsep=0.3em]
\item[Arbitrage pur.] Stratégie sans coût initial, sans possibilité de perte, avec une possibilité de gain (chapitre~1).
\item[Arbitrage statistique.] Stratégie sans coût initial dont le gain espéré est positif et dont la probabilité de perte tend vers zéro à long terme (chapitre~1).
\item[Base.] Écart entre le prix à terme et le prix comptant d'un même sous-jacent (chapitre~14).
\item[\emph{Cash-and-carry}.] Achat du comptant financé et vente du contrat à terme, pour capter la base (chapitre~14).
\item[Cointégration.] Relation entre séries $I(1)$ dont une combinaison linéaire est stationnaire (chapitre~4).
\item[Copule.] Fonction qui décrit la dépendance entre variables indépendamment de leurs lois marginales (chapitre~9).
\item[CUSUM.] Somme cumulée des écarts standardisés, utilisée pour détecter un changement de régime (chapitre~17).
\item[Demi-vie.] Durée au bout de laquelle l'écart espéré à la moyenne est divisé par deux (chapitres~1 et~5).
\item[\emph{Drawdown}.] Perte depuis le plus haut historique de la courbe de gains.
\item[\emph{Funding}.] Paiement périodique entre acheteurs et vendeurs d'un contrat perpétuel, qui ramène son prix vers l'indice (chapitre~14).
\item[Jambe.] Chacun des instruments d'une position d'arbitrage. \emph{Risque de jambe} : une jambe exécutée, l'autre non ou en retard (chapitre~15).
\item[\emph{Maker}, \emph{taker}.] Ordre qui apporte de la liquidité au carnet (limite, en attente), ordre qui en prend (exécuté immédiatement) (chapitre~12).
\item[\emph{Paper trading}.] Exécution simulée en temps réel, sans argent engagé (chapitre~18).
\item[Parité.] Égalité des décisions entre la version de recherche et la version de production d'une stratégie (chapitre~18).
\item[Point de base (pb).] Un centième de pourcent : $10^{-4}$.
\item[\emph{Prop firm}.] Société qui finance des traders sous conditions ; ses règles interdisent en général l'arbitrage de latence et la couverture croisée entre comptes.
\item[Régression fallacieuse.] Relation apparente entre séries non stationnaires indépendantes (chapitre~1).
\item[Sélection adverse.] Tendance d'un ordre \emph{maker} à être exécuté précisément quand le marché lui donne tort (chapitre~15).
\item[Stationnarité.] Moyenne, variance et autocovariances constantes dans le temps (chapitre~1).
\item[Surapprentissage.] Choix de paramètres qui décrivent le bruit de l'échantillon plutôt qu'une régularité durable (chapitre~16).
\item[\emph{Walk-forward}.] Validation glissante : estimer sur le passé, trader sur la période suivante, avancer (chapitre~16).
\item[\emph{Z-score}.] Écart standardisé $(S_t - m)/d$ (chapitre~1).
\end{description}

% =====================================================================
\chapter{Reproduire les résultats du cours}
\label{ann:repro}
% =====================================================================

\section*{Installation}
\begin{lstlisting}[language=bash]
git clone https://github.com/OptiWealthSolutions/arbitrage_hub.git
cd arbitrage_hub
python -m venv .venv && source .venv/bin/activate
pip install -r requirements.txt
pytest -q                                   # 24 tests
\end{lstlisting}

\section*{Scripts par chapitre}

Chaque chiffre du cours provient d'un script à graine fixée ; relancer le script redonne exactement les mêmes nombres. Les résultats sont écrits dans \texttt{code/chXX/results.json} (\texttt{code/ch01/results/} pour le chapitre~1) et les figures dans \texttt{docs/cours/figures/chXX/}.

\begin{center}\small
\begin{longtable}{>{\raggedright}p{1.1cm}>{\ttfamily\footnotesize}p{6.4cm}>{\raggedright\arraybackslash}p{6.0cm}}
\toprule
Ch. & \normalfont Script & Contenu \\
\midrule
\endhead
1 & ch01/h1 à h5\_*.py & marche aléatoire, AR(1) et demi-vie, corrélation contre cointégration, régression fallacieuse, règle en \emph{z-score} \\
2 & ch02/ch02\_regression.py & MCO, asymétrie, régression orthogonale, vraisemblance AR(1), biais \\
3 & ch03/ch03\_stationarity\_tests.py & Dickey-Fuller, puissance, KPSS, ratio de variance, Hurst, cas réels \\
4 & ch04/ch04\_cointegration.py & Engle-Granger, correction d'erreur, Johansen, cointégration partielle \\
5 & ch05/ch05\_ornstein\_uhlenbeck.py & simulation exacte, estimation, temps de passage, seuils optimaux \\
6 & ch06/ch06\_distance\_method.py & méthode des distances (Gatev et al.) \\
7 & ch07/ch07\_cointegration\_selection.py & sélection de paires, faux positifs, Bonferroni, BH \\
8 & ch08/ch08\_kalman.py & filtre de Kalman, ratio variable \\
9 & ch09/ch09\_copulas.py & copules, fonction $h$, cas KO/PEP \\
10 & ch10/ch10\_pca\_statarb.py & ACP, Marchenko-Pastur, \emph{s-score} \\
11 & ch11/ch11\_machine\_learning.py & fuite d'information, validation purgée, sélection par \emph{clustering} \\
12 & ch12/ch12\_microstructure.py & fourchette, Roll, impact, coûts d'un aller-retour \\
13 & ch13/ch13\_spatial\_triangular.py & arbitrage spatial, triangulaire, course de latence \\
14 & ch14/ch14\_basis\_funding.py & portage, \emph{cash-and-carry}, \emph{funding}, liquidation \\
15 & ch15/ch15\_execution.py & risque de jambe, sélection adverse \\
16 & ch16/ch16\_backtest.py & biais d'anticipation, Sharpe, DSR, PBO, survie \\
17 & ch17/ch17\_risk.py & Kelly, protections contre la rupture, spirale de marge, diversification \\
18 & ch18/ch18\_pipeline.py & parité, critères de passage, retard d'exécution, suivi du \emph{paper trading} \\
\bottomrule
\end{longtable}
\end{center}

\section*{Recompiler ce document}
\begin{lstlisting}[language=bash]
cd docs/cours && sh build.sh     # formulaire regénéré, 3 passes LaTeX + BibTeX
\end{lstlisting}

\section*{Données réelles}
Les études de cas utilisent 500 clôtures quotidiennes de GLD, GLDM, KO et PEP (21 octobre 2024 au 19 octobre 2026) téléchargées via Interactive Brokers ; voir \texttt{data/case\_studies/README.md}. Prix non ajustés des dividendes.

\section*{Avertissement}
Ce cours est un support pédagogique et non un conseil en investissement. Les performances citées, qu'elles proviennent de la littérature ou des simulations, ne préjugent pas des performances futures. Vérifiez les conditions de vos courtiers et plateformes avant tout déploiement ; les stratégies d'arbitrage inter-plateformes ne doivent pas être utilisées sur des comptes de \emph{prop firms}.
